\chapter{Just because you decide one day to take up this most unusual career (I Think That Would Be Truly Wonderful)}

The magical girl transformation sequence has a curious history and function, which makes the recapitulation of Mami's transformation a fitting opening for Madoka Magica's second episode. The sequence originates in Go Nagai's 1973 manga and anime Cutie Honey, itself interesting as a hybrid of magical girl and shonen (that is, "for boys") elements \cite{ref24}. This sequence soon becomes an ubiquitous feature of the genre, and by the 1990s magical girl shows typically hybridize shonen elements \cite{ref25}. 

Within most magical girl shows, the transformation sequence serves a number of functions. The traditional sequence (epitomized by Sailor Moon) involves colored silhouettes of apparently nude characters, on whom the magical girl costume forms. The magical girl frequently takes a sequence of poses as the camera spins around her, creating the effect of a lone nude girl dancing in the center of a multitude of Male Gazes \cite{ref26}. In story terms, this is an empowering moment, and the music that accompanies it usually evokes a similar feeling, but the camera simultaneously asserts the character as an object to be looked at; she is, in other words, empowered, but not too empowered. She is still performing femininity and submitting herself to masculine hegemony \cite{ref27}. The sequence thus serves to say, more or less, "Yes, this young woman is a powerful protector, but she is still not going to threaten the patriarchal status quo or your fragile male ego; she is here for your enjoyment." There is a reason the magical girl genre is increasingly associated with the stereotypical basement-dwelling pervert \cite{ref28}. 

In his Beautiful Fighting Girl, Japanese psychiatrist Tamaki Saito defines the otaku-broadly, Japanese for "geek, nerd, fan," albeit with somewhat harsher connotations than the English equivalent-as someone who, among other criteria, possesses the ability to transform the fiction they consume into something they feel they own as a part of their identity, as well as being able to find sexual objects in fiction. In other words, a key element of otaku culture in Japan-one which can be observed in American geek culture as well-is this idea that the work in question and its characters belong to the viewer, and this is intimately entwined with the sexual desires of the (typically heterosexual male) viewer. In that light, the transformation sequence represents, as Saito notes, an accelerated maturation, which is to say a sort of tacit permission by the show for viewers to sexualize its characters. In Saito's eyes, the magical girl (or "fighting girl" as he calls it, extending the concept to cover characters in science fiction and other genres) is neither an empowering figure permitting an audience of girls to feel strong, nor a tool for teaching girls traditional feminine virtues; she is not someone for the viewer to become but rather someone for the viewer to possess, generally sexually \cite{ref29}. 

Mami's sequence is tamer than most, seeing as she never loses her school uniform, but still lingers on breast and buttocks, emphasizing her body and costume rather than her face. Mami is a feminine figure, maternal and kind even as she quite violently blows away the familiar threatening Madoka and Mami. This transformation is paralleled only a few minutes later by Junko, who starts her only scene of the episode in a mothering role, gently, kindly chiding Madoka for staying out late the previous night. As soon as she finishes putting her makeup on, however, she becomes the steely, ambitious executive, weighing her allies and options in a potential bid for dominance of her company. In both cases, the transformation sequence serves as a gateway from the traditionally feminine role of schoolgirl or mother to a traditionally masculine role as warrior or conqueror.

For Mami in particular, the transformation sequence serves as a way to interrupt the strangeness of the witch's labyrinth with a normative element of magical girl shows that reinforces both the norms of the genre and the social norms of patriarchal society. It may also be a performance for the benefit of Madoka and Sayaka, a familiar signifier that Mami is a heroic figure come to rescue them. Notably, however, when she transforms again near the end of the episode there is no such sequence, just a quick burst of light and ribbons after which Mami is in her magical girl costume. Where the full transformation sequence emphasizes that the schoolgirl exists within the magical girl, who exists for the consumption of the viewer (Madoka and Sayaka diegetically, and the Male Gaze extradiegetically), this abbreviated sequence does the opposite, reminding the viewer that the warrior exists within the woman and denying the viewer the opportunity to possess and consume her. Rather than reinforcing the social order, it undermines it, creating the suggestion that maybe the magical girl doesn't need to perform for a masculine audience or emphasize her femininity to gain their permission to fight; perhaps all she needs is a cause.

Certainly that seems to be enough for Sayaka. She is drawn in immediately by Mami's role as protector of the innocent, willing immediately to cheer for Mami and hate Homura. For Sayaka, the magical girl represents an exciting conflict between a clear good and a clear evil, one Sayaka is eager to join on the side of good. She wants to become a magical girl so she can fight for justice and stand by Mami's side, hence her decision to bring a weapon to the witch-hunting session. By contrast, Madoka is more interested in the experience of being a magical girl; she wants to understand Mami and Homura both, to make friends with them. She wants to become a magical girl so that she can feel "cool" and special, and to be like Mami, hence her focus on designing a costume. Put another way, Sayaka is focused on doing and Madoka is focused on being. Sayaka is active, but unfocused; Madoka is centered, but passive.

Regardless of their different interests in becoming a magical girl, both girls struggle with the question of the wish. Surprisingly for such young girls, the two meditate briefly on how privileged their lives are; as children, they have limited influence on the quality of their own lives, and so must acknowledge that the fact that they are safe, healthy, fairly well-educated for their age, members of their society's dominant ethnicity, good-looking, and well-off (the apartment-dwelling Sayaka appears to be somewhat less wealthy than Madoka, whose house is enormous, but based on later interactions with Kyoko, neither appears to have experience with food or housing insecurity) is down entirely to their luck in being born to their particular parents, not any special effort on their part. It is Sayaka who articulates this problem; as we will see later, she has seen in Kyousuke how easily random events can derail a life. However, Madoka shares her sentiments; both girls feel that they are simply too well-off to have any wishes worth the price Kyubey charges.

There is an underlying theme throughout this episode that posits magical girls as part of an ecosystem of sorts. Kyoko will state this outright in a few episodes, but already we have the basics: witches predate on humans, and magical girls predate on witches. Within this ecosystem, Kyubey is a scavenger, feeding on the byproducts of the predation, but socially he is pure predator, picking out the most vulnerable members of the group to target and use for his purposes. Based on Mami's, Homura's, and Kyoko's flashbacks to the circumstances of their contracts, it is clear that he seeks out girls who are in pain and despair, and thus likely both to be willing to take the contract and to turn relatively quickly into witches. Madoka and Sayaka's slowness to accept his offer is a product of their fortunate circumstances; it is quite difficult to manipulate someone who doesn't need anything.

This is usually not an issue with magical girls. Kyubey seeks consent from the girls who fight his war, and within the context of these first two episodes thus comes off looking rather better than his equivalents in Cardcaptor Sakura and Sailor Moon. Neither Sakura nor Serena was offered a choice in their respective shows; the former was told that as the one who released the Clow, it was her duty to recover them \cite{ref30}, and the latter simply informed that she was destined to become Sailor Moon. As is often the case in both media and real life, women are expected to simply accept that they will fulfill certain roles, regardless of their consent \cite{ref31}. (Which is not to say that this doesn't happen to men, just that men in media are more likely to be offered a choice and more likely to resist their fight. Shinji in Neon Genesis Evangelion struggles constantly against his destiny as an Eva pilot; Serena/Usagi just goes along with being a magical girl.)

Knowledge of later episodes, of course, makes Kyubey's duplicity obvious. He is not seeking informed consent at all, but rather tricking and manipulating women in truly appalling ways. Even within these initial episodes, familiarity with stories of trickster genies and deals with the devil (emphasized by the graffiti quoting Goethe's Faust that the witch victim walks past) suggests that Kyubey is not what he seems and not to be trusted. Kyubey's unreadable facial expression adds to this discomfort; it may initially seem cute or friendly, but as the viewer grows more suspicious of him, Kyubey's lack of affect begins to seem more like a mask that makes him harder to read, compounding the suspicion. Madoka and Sayaka are not living in Cardcaptor Sakura (no matter how much Madoka's waking at the beginning of the episode mirrors the second episode of that show \cite{ref32}) or Sailor Moon, and the world of the magical girl will soon seep into and corrupt their safe, comfortable, privileged lives whether they make the deal with Kyubey or not. 

Mami may seek to protect the innocent and maintain the norms of the magical girl genre, but she will fail. Though she transforms Sayaka's bat into a pink, cartoonish scepter of the sort that might appear in the typical magical girl show, the bat is useless against the witch at the end of the episode, and serves only as a means to temporarily lock Sayaka and Madoka away from it while Mami fights. "Magia," the harbinger of the dark version of Madoka Magica, plays throughout this battle, as the witch imposes its surreal labyrinth onto Madoka's world. Mami is able to defeat it once again, but it puts up an admirable fight that reveals both her overconfidence (as she falls for the butterfly rope trap) and constant awareness of her audience, for whom she is clearly performing rather than focusing on the immediate threat of the witch.

Mami, by the end of this episode, has firmly positioned herself as a protector and mentor to Sayaka and Madoka. So long as she is around, she will keep the "Magia" version of Madoka Magica confined and at bay, preventing it from overwhelming what might be called the "Mata Ashita" version of the show; that is, she will keep the emotional intensity, complexity, and menace of the witches back and maintain the stability and comfort of Madoka's happy, safe-for-preteen-viewers little life. However, this episode also establishes her weakness. It is no longer merely clear that Mami must go for the show to move on; it is clear how she will go, as a victim of her own confidence and tendency to showboat.
